\chapter{Desarrollos matemáticos complementarios}
\label{ap:desarrollos}

Este apéndice recoge dos derivaciones autocontenidas a las que se remite desde
el cuerpo del trabajo: la forma general de las condiciones de Karush-Kuhn-Tucker
(Capítulo~\ref{ch:formulacion}) y la construcción del estimador de covarianza
por componentes principales junto con su equivalencia con el modelo APT
(Capítulo~\ref{ch:limitaciones}).

\section{Condiciones de Karush-Kuhn-Tucker: planteamiento general}
\label{ap:kkt-general}

Consideremos el problema de optimización con restricciones de igualdad
y desigualdad:
\begin{equation}
    \begin{aligned}
        \min_{\mathbf{x}} \quad & f(\mathbf{x}) \\
        \text{s.a.} \quad & g_i(\mathbf{x}) \le 0, \quad i = 1, \dots, m, \\
        & h_j(\mathbf{x}) = 0, \quad j = 1, \dots, p.
    \end{aligned}
    \label{eq:problema-general}
\end{equation}
El \textbf{Lagrangiano} asociado es:
\begin{equation}
    \mathcal{L}(\mathbf{x}, \boldsymbol{\lambda}, \boldsymbol{\nu})
    = f(\mathbf{x})
    + \sum_{i=1}^{m} \lambda_i g_i(\mathbf{x})
    + \sum_{j=1}^{p} \nu_j h_j(\mathbf{x}),
    \label{eq:lagrangiano-general}
\end{equation}
donde $\lambda_i \ge 0$ son los multiplicadores de las restricciones de
desigualdad y $\nu_j \in \mathbb{R}$ los de las igualdades. En el
Capítulo~\ref{ch:formulacion} este planteamiento se particulariza al programa
cuadrático de Markowitz.

\section{Lagrangiano del QP de Markowitz}
\label{ap:kkt-markowitz}

Escribiendo el QP~(\ref{eq:qp-completo}) en la forma estándar anterior, la
función objetivo es $f(\mathbf{w}) = \frac{1}{2}\mathbf{w}^T \boldsymbol{\Sigma}
\mathbf{w}$ y las restricciones de desigualdad $g_i(\mathbf{w}) \le 0$ son:
\begin{align}
    g_i^{(1)}(\mathbf{w}) &= -w_i \le 0, \quad i = 1, \dots, n,
        &&\text{(long-only)}, \label{eq:g1} \\
    g_i^{(2)}(\mathbf{w}) &= w_i - w_{\max} \le 0, \quad i = 1, \dots, n,
        &&\text{(tope por activo)}, \label{eq:g2} \\
    g^{(3)}(\mathbf{w}) &= \mu^* - \mathbf{w}^T \boldsymbol{\mu} \le 0,
        &&\text{(rentabilidad objetivo)}, \label{eq:g3}
\end{align}
junto con la restricción de igualdad
\begin{equation}
    h(\mathbf{w}) = \mathbf{1}^T \mathbf{w} - 1 = 0. \label{eq:h}
\end{equation}
Introduciendo los multiplicadores $\alpha_i \ge 0$ (cota inferior),
$\beta_i \ge 0$ (cota superior), $\gamma \ge 0$ (rentabilidad objetivo) y
$\lambda \in \mathbb{R}$ (presupuesto), el Lagrangiano es:
\begin{equation}
    \begin{aligned}
        \mathcal{L}(\mathbf{w}, \boldsymbol{\alpha}, \boldsymbol{\beta},
        \gamma, \lambda)
        =& \; \frac{1}{2} \mathbf{w}^T \boldsymbol{\Sigma} \mathbf{w}
           - \sum_{i=1}^{n} \alpha_i w_i
           + \sum_{i=1}^{n} \beta_i (w_i - w_{\max}) \\
        & + \gamma (\mu^* - \mathbf{w}^T \boldsymbol{\mu})
           + \lambda (1 - \mathbf{1}^T \mathbf{w}).
    \end{aligned}
    \label{eq:lagrangiano-qp}
\end{equation}
Las condiciones KKT que de aquí se derivan se enuncian e interpretan en el
Capítulo~\ref{ch:formulacion}.

\section{Estimador de covarianza por PCA y equivalencia con la APT}
\label{ap:pca-apt}

La matriz de covarianza muestral admite una descomposición espectral:
\begin{equation}
    \boldsymbol{\Sigma} = \mathbf{V} \boldsymbol{\Lambda} \mathbf{V}^T
    = \sum_{i=1}^{n} \lambda_i \mathbf{v}_i \mathbf{v}_i^T,
    \label{eq:descomposicion-espectral}
\end{equation}
donde $\lambda_1 \ge \lambda_2 \ge \dots \ge \lambda_n \ge 0$ son los
autovalores (ordenados de mayor a menor) y $\mathbf{v}_1, \dots, \mathbf{v}_n$
los autovectores. Cada autovalor $\lambda_i$ representa la varianza explicada
por el $i$-ésimo componente principal. El estimador PCA retiene únicamente los
$k$ primeros componentes (los factores principales) y descarta los $n-k$
restantes, que se atribuyen a ruido:
\begin{equation}
    \boldsymbol{\Sigma}_k = \sum_{i=1}^{k} \lambda_i \mathbf{v}_i
    \mathbf{v}_i^T,
    \label{eq:sigma-k}
\end{equation}
y conserva la varianza específica de cada activo no explicada por los factores
comunes (la diagonal del residuo):
\begin{equation}
    \boldsymbol{\Sigma}_{\text{PCA}} = \boldsymbol{\Sigma}_k
    + \mathop{\textnormal{diag}}\nolimits\bigl(\boldsymbol{\Sigma}
    - \boldsymbol{\Sigma}_k\bigr).
    \label{eq:sigma-pca-final}
\end{equation}

La Arbitrage Pricing Theory de Ross \cite{Ross1976} postula que las
rentabilidades se generan por un modelo factorial:
\begin{equation}
    r_i = \alpha_i + \sum_{j=1}^{k} \beta_{ij} f_j + \varepsilon_i,
    \label{eq:apt}
\end{equation}
donde $f_1, \dots, f_k$ son factores de riesgo sistemático y $\varepsilon_i$ es
el riesgo idiosincrásico del activo $i$, no correlacionado con los factores ni
con los $\varepsilon$ de otros activos. Bajo este modelo, la covarianza se
descompone como
\begin{equation}
    \boldsymbol{\Sigma} = \mathbf{B} \boldsymbol{\Sigma}_f
    \mathbf{B}^T + \boldsymbol{\Psi},
    \label{eq:cov-apt}
\end{equation}
con $\mathbf{B}$ la matriz de cargas factoriales ($\beta_{ij}$),
$\boldsymbol{\Sigma}_f$ la covarianza de los factores y $\boldsymbol{\Psi}$ una
matriz diagonal con las varianzas idiosincrásicas. Si se toman como factores los
$k$ primeros componentes principales de las rentabilidades, el
estimador~(\ref{eq:sigma-pca-final}) es exactamente el estimador APT: las cargas
vienen dadas por $\mathbf{B} = \mathbf{V}_k \boldsymbol{\Lambda}_k^{1/2}$ y la
covarianza de los factores es $\boldsymbol{\Sigma}_f = \mathbf{I}_k$ (por
ortogonalidad de los componentes principales). Esta equivalencia conecta el
enfoque puramente estadístico (PCA) con la teoría financiera (APT).
